% \documentclass[a4paper,12pt]{article}
\documentclass[a4paper,11pt]{article}
% -------------------------------------------------
% Set up the page margins.
% -------------------------------------------------
% \usepackage[text={6.5in,9.5in},centering]{geometry}
% \setlength{\topmargin}{-0.25in}
\usepackage[margin=1in]{geometry}

% -------------------------------------------------
% Load Packages here
% -------------------------------------------------
%\usepackage{hyperref}
\usepackage{graphicx,url,subfig}
\RequirePackage[OT1]{fontenc}
\RequirePackage{amsthm,amsmath,amssymb,amscd}
\RequirePackage[numbers]{natbib}
\RequirePackage[colorlinks,citecolor=blue,urlcolor=blue]{hyperref}
\usepackage{graphics}
\usepackage{enumerate}
\usepackage{datetime}
\usepackage{etex}
% \usepackage[notcite,notref,color]{showkeys}
\usepackage[normalem]{ulem} % either use this (simple) or
\usepackage{soul} % use this (many fancier options)
\makeatother
\numberwithin{equation}{section}
\allowdisplaybreaks
\usepackage{mathrsfs}


% -------------------------------------------------
% check references
% -------------------------------------------------
% \usepackage{refcheck}
% \usepackage[notref,notcite]{showkeys}
% \usepackage[notcite]{showkeys}
% \usepackage{showkeys}

% -------------------------------------------------
% Environments here
% -------------------------------------------------
\theoremstyle{plain}
\newtheorem{theorem}{Theorem}[section]
\newtheorem{corollary}[theorem]{Corollary}
\newtheorem{lemma}[theorem]{Lemma}
\newtheorem{proposition}[theorem]{Proposition}
\newtheorem{exercise}{Exercise}
\theoremstyle{definition}
\newtheorem{definition}[theorem]{Definition}
\newtheorem{hypothesis}[theorem]{Hypothesis}
\newtheorem{assumption}[theorem]{Assumption}

\newtheorem{remark}[theorem]{Remark}
\newtheorem{conjecture}[theorem]{Conjecture}
\newtheorem{example}[theorem]{Example}

% -------------------------------------------------
%  Define short hand symbols.
% -------------------------------------------------
\newcommand{\E}{\mathbb{E}}
\newcommand{\W}{\dot{W}}
\newcommand{\B}{\textbf{B}}
\newcommand{\D}{\mathbb{D}}
\newcommand{\ud}{\ensuremath{\mathrm{d}}}
\newcommand{\Ceil}[1]{\left\lceil #1 \right\rceil}
\newcommand{\Floor}[1]{\left\lfloor #1 \right\rfloor}
\newcommand{\sgn}{\text{sgn}}
\newcommand{\Lad}{\text{L}_{\text{ad}}^2}
\newcommand{\SI}[1]{\mathcal{I}\left[#1 \right]}
\newcommand{\SIB}[2]{\mathcal{I}_{#2}\left[#1 \right]}
\newcommand{\Indt}[1]{1_{\left\{#1 \right\}}}
\newcommand{\LadInPrd}[1]{\left\langle #1 \right\rangle_{\text{L}_\text{ad}^2}}
\newcommand{\LadNorm}[1]{\left|\left|  #1 \right|\right|_{\text{L}_\text{ad}^2}}
\newcommand{\Norm}[1]{\left\|  #1   \right\|}
\newcommand{\Ito}{It\^{o} }
\newcommand{\Itos}{It\^{o}'s }
\newcommand{\spt}[1]{\text{supp}\left(#1\right)}
\newcommand{\InPrd}[1]{\left\langle #1 \right\rangle}
\newcommand{\mr}{\textbf{r}}
\newcommand{\Ei}{\text{Ei}}
\newcommand{\arctanh}{\operatorname{arctanh}}
\newcommand{\ind}[1]{\mathbb{I}_{\left\{ {#1} \right\} }}

\DeclareMathOperator{\esssup}{\ensuremath{ess\,sup}}

\newcommand{\steps}[1]{\vskip 0.3cm \textbf{#1}}
\newcommand{\calB}{\mathcal{B}}
\newcommand{\calD}{\mathcal{D}}
\newcommand{\calE}{\mathcal{E}}
\newcommand{\calF}{\mathcal{F}}
\newcommand{\calG}{\mathcal{G}}
\newcommand{\calK}{\mathcal{K}}
\newcommand{\calH}{\mathcal{H}}
\newcommand{\calI}{\mathcal{I}}
\newcommand{\calL}{\mathcal{L}}
\newcommand{\calM}{\mathcal{M}}
\newcommand{\calN}{\mathcal{N}}
\newcommand{\calO}{\mathcal{O}}
\newcommand{\calT}{\mathcal{T}}
\newcommand{\calP}{\mathcal{P}}
\newcommand{\calR}{\mathcal{R}}
\newcommand{\calS}{\mathcal{S}}
\newcommand{\bbC}{\mathbb{C}}
\newcommand{\bbN}{\mathbb{N}}
\newcommand{\bbP}{\mathbb{P}}
\newcommand{\bbZ}{\mathbb{Z}}
\newcommand{\myVec}[1]{\overrightarrow{#1}}
\newcommand{\sincos}{\begin{array}{c} \cos \\ \sin \end{array}\!\!}
\newcommand{\CvBc}[1]{\left\{\:#1\:\right\}}
\newcommand*{\one}{{{\rm 1\mkern-1.5mu}\!{\rm I}}}
\newcommand{\RR}{\mathbb{R}}
\newcommand{\R}{\mathbb{R}}
\newcommand{\myEnd}{\hfill$\square$}
\newcommand{\ds}{\displaystyle}
\newcommand{\Shi}{\text{Shi}}
\newcommand{\Chi}{\text{Chi}}
\newcommand{\Daw}{\ensuremath{\mathrm{daw}}}
\newcommand{\Erf}{\ensuremath{\mathrm{erf}}}
\newcommand{\Erfi}{\ensuremath{\mathrm{erfi}}}
\newcommand{\Erfc}{\ensuremath{\mathrm{erfc}}}
\newcommand{\He}{\ensuremath{\mathrm{He}}}

\def\crtL{\theta}

\newcommand{\conRef}[1]{(#1)}

\DeclareMathOperator{\Lip}{\mathit{L}}
\DeclareMathOperator{\LIP}{Lip}
\DeclareMathOperator{\lip}{\mathit{l}}

\DeclareMathOperator{\Vip}{\overline{\varsigma}}
\DeclareMathOperator{\vip}{\underline{\varsigma}}
\DeclareMathOperator{\vv}{\varsigma}

% % Set up mdframe for questions.
\newcounter{NQ}
\newcounter{LQ}
\newcounter{ZQ}
\usepackage{xcolor}
\usepackage[framemethod=tikz]{mdframed}
\mdfsetup{%
	middlelinecolor=lightgray,
	middlelinewidth=2pt,
	backgroundcolor=red!20,
	roundcorner=10pt}
\newenvironment{NickQuestion}[1][] %
{\refstepcounter{NQ}
	\begin{mdframed}[backgroundcolor=lightgray]\textbf{Nick's Question \theNQ \: #1~~}~~~\noindent}%
	{\end{mdframed}}
\newenvironment{LeQuestion}[1][] %
{\refstepcounter{LQ}
	\begin{mdframed}[backgroundcolor=red!10]\textbf{Le's Question \theLQ \: #1~~}~~~\noindent}%
	{\end{mdframed}}
\newenvironment{ZhijianQuestion}[1][] %
{\refstepcounter{ZQ}
	\begin{mdframed}[backgroundcolor=blue!10]\textbf{Zhijian's Question \theZQ \: #1~~}~~~\noindent}%
	{\end{mdframed}}
\numberwithin{NQ}{section}
\numberwithin{LQ}{section}
\numberwithin{ZQ}{section}

\usepackage[notref,notcite]{showkeys}

\title{Nick Eisenberg's Meeting Notes}
\author{Le Chen, Nick Eisenberg, Zhijian Wu}
\date{October 2019}

\usepackage{natbib}
\usepackage{graphicx}


\begin{document}

\begin{NickQuestion}
In the Wave equation paper you showed that for the Greens function, we have that 
\[
	\mathcal{F}G(t,\cdot)(\xi)=\frac{\sin(t |\xi|)}{|\xi|}
\]
and so we see that
\[
	\mathcal{F}G(t,\cdot)(c\xi) = \frac{1}{c}\mathcal{F}(ct, \cdot)(\xi)
\]

However in the Fractional SPDE paper the greens function is more complicated and we have that 

\begin{align*}
	\mathcal{F}Y_{\alpha, \beta, \gamma, d}(t, \cdot)(\xi) &= t^{\beta + \gamma -1} E_{\beta, \beta + \gamma}(-2^{-1}\nu t^\beta |\xi|^\alpha)
	\\&= t^{\beta + \gamma -1} H^{1,1}_{1,2}\left(2^{-1}\nu t^\beta |\xi|^\alpha \bigg\vert \begin{matrix} (0,1) &
	\\ (0,1) & (1- \beta - \gamma ,\beta) \end{matrix} \right)
	\\&= t^{\beta + \gamma -1} \frac{1}{2 \pi i} \int_L \mathcal{H}^{1\;1}_{1\;2}(s) \left(2^{-1}\nu t^\beta |\xi|^\alpha\right)^{-s} \ud s
\end{align*}

And so we see that 
\[
\mathcal{F}Y_{\alpha, \beta, \gamma, d}(t, \cdot)(c\xi) =	t^{\beta + \gamma -1} \frac{1}{2 \pi i} \int_L \mathcal{H}^{1\;1}_{1\;2}(s) \left(2^{-1}\nu t^\beta |c\xi|^\alpha\right)^{-s} \ud s
\]

However I am unsure how to handle the constant, $c$, on the inside so we can get a nice scaling property like we do in the Wave Equation paper. 

Also, instead of using the Fox-H function we can just say 
\begin{align*}
		\mathcal{F}Y_{\alpha, \beta, \gamma, d}(t, \cdot)(\xi) &= t^{\beta + \gamma -1} E_{\beta, \beta + \gamma}(-2^{-1}\nu t^\beta |\xi|^\alpha)
		\\&= t^{\beta + \gamma -1} \sum_{k=0}^\infty \dfrac{\left(-2^{-1}\nu t^\beta |\xi|^\alpha\right)^k}{\Gamma\left(\beta k + (\beta + \gamma)\right)}
\end{align*}
and so 
\[
		\mathcal{F}Y_{\alpha, \beta, \gamma, d}(t, \cdot)(c\xi) = t^{\beta + \gamma -1} \sum_{k=0}^\infty \dfrac{\left(-2^{-1}\nu t^\beta |c\xi|^\alpha\right)^k}{\Gamma\left(\beta k + (\beta + \gamma)\right)}
\]
but I still don't see how to handle the $c$. 
\end{NickQuestion}

\bigskip
\textcolor{magenta}{
Here is how one should proceed:
\begin{align*}
    \mathcal{F}Y_{\alpha, \beta, \gamma, d}(t, \cdot)(c\xi) & = t^{\beta + \gamma -1} \sum_{k=0}^\infty \dfrac{\left(-2^{-1}\nu t^\beta |c\xi|^\alpha\right)^k}{\Gamma\left(\beta k + (\beta + \gamma)\right)} \\
                                                            & = t^{\beta + \gamma -1} \sum_{k=0}^\infty \dfrac{\left(-2^{-1}\nu t^\beta c^\alpha |\xi|^\alpha\right)^k}{\Gamma\left(\beta k + (\beta + \gamma)\right)}\\
							    & = t^{\beta + \gamma -1} \sum_{k=0}^\infty \dfrac{\left(-2^{-1}\nu (tc^{\frac{\alpha}{\beta}})^\beta |\xi|^\alpha\right)^k}{\Gamma\left(\beta k + (\beta + \gamma)\right)}\\
							    & =c^{\frac{\alpha}{\beta}(1-\gamma-\beta)} (tc^{\frac{\alpha}{\beta}})^{\beta + \gamma -1} \sum_{k=0}^\infty \dfrac{\left(-2^{-1}\nu (tc^{\frac{\alpha}{\beta}})^\beta |\xi|^\alpha\right)^k}{\Gamma\left(\beta k + (\beta + \gamma)\right)}\\
							    & =c^{\frac{\alpha}{\beta}(1-\gamma-\beta)} \mathcal{F}Y_{\alpha, \beta, \gamma, d}\left(c^{\frac{\alpha}{\beta}}t, \cdot\right)(c\xi) 
\end{align*}
Please check if this scaling reduces to heat and wave equations correctly.}


\begin{NickQuestion}
	In the wave equation paper with Xia Chen you said that that fundamental solution to the SPDE
\begin{equation}
\begin{cases}
\dfrac{\partial^2 u}{\partial t^2}(t,x)-\frac{1}{2}\Delta u(t,x) = \sqrt{\Theta}u(t,x)\dot W(x) & \text{$x\in \R^d$ , $t>0$} \\
u(0,\cdot) = 1  & 
\\ u_t(0,\cdot) = 0
\end{cases}  \label{wspde}
\end{equation} 
and the mild solution is defined to be 
\begin{equation} \label{weqs}
u(t,x) = 1 + \sqrt{\Theta}\int_0^t \left( \int_{\R^d}G(t-s,x-y)u(s,y) W(\delta y \right) \ud s 
\end{equation}
but for general initial conditions for the wave equation we have the solution to be 

\begin{align}u(t,x) = \dfrac{\partial}{\partial t} \int_{\R^d}G(t-s,x-y)u(0,y) \ud y &+ \int_{\R^d}G(t-s,x-y)u_t(0,y) \ud y \nonumber
\\& + \sqrt{\Theta}\int_0^t \left( \int_{\R^d}G(t-s,x-y)u(s,y) W(\delta y) \right) \ud s \label{gweqs}
\end{align}

But when the initial conditions are such as in \ref{wspde} then the first two integral in \ref{gweqs} are $1$ and $0$ respectively and so we arrive at \ref{weqs}.

In the fractional SPDE
\begin{equation}
\begin{cases}
\left(\partial^\beta + \frac{\nu}{2}(-\Delta)^{\alpha/2}\right)u(t,x) = I^\gamma_t \left[\rho(u(t,x))\dot W(x) \right]& \text{$x\in \R^d$ , $t>0$} \\
u(0,\cdot) = \mu  & \beta \in (0,1]
\\ u(0,\cdot) = \mu_0, \quad u_t(0,\cdot) = \mu_1 & \beta \in (1,2]
\end{cases}  \label{fspde}
\end{equation} 
and we define 
\[
J_0(t.x) = \begin{cases}
						(Z(t,\cdot)*\mu)(x) & \beta \in (0,1]
						\\ 	(Z^*(t,\cdot)*\mu_0)(x) +  (Z(t,\cdot)*\mu_1)(x) & \beta \in (0,1]
				\end{cases}	
\]
and
\[			
I(t,x) \iint_{[0,t] \times \R} Y(t-s, x-y) \rho (u(s,y)) W(\delta y) \ud s
\]
and the general solution for \ref{fspde} is 
\begin{align}
u(t,x) = J_0(t,x) + I(t,x) \label{feqs}
\end{align}

My question is that if we consider the following fractional spde
 \begin{equation}
 \begin{cases}
 \left(\partial^\beta + \frac{\nu}{2}(-\Delta)^{\alpha/2}\right)u(t,x) = I^\gamma_t \left[\rho(u(t,x))\dot W(x) \right]& \text{$x\in \R^d$ , $t>0$} \\
 u(0,\cdot) = 1  & \beta \in (0,1]
 \\ u(0,\cdot) = 1, \quad u_t(0,\cdot) = 0 & \beta \in (1,2]
 \end{cases}  \label{fspde1}
 \end{equation} 
 then can we calculate $J_0(t,x)$? More specifically is $J_0(t,x) = 1$ in this case? The wave equation paper worked out nicely because in that paper $J_0(t,x) =1$ and so the solution was able to be written as \ref{weqs}. The $1$ allows us to easily do the Piccard iteration to get equation (2.3) from the Wave equation paper with Xia Chen. 
 
 So if our solution can not be written as 
 \[
 u(t,x) = 1 + I(t,x)
 \] 
 I am not sure at the moment how we can come up with an expression for the $f_n's$ in the Weiner Expansion like you did in equation (2.3) in the paper with Xia Chen. 
\end{NickQuestion}
\end{document}
